% ============================================================
%  ChemoInteract Documentation - Part 2
%  Title: Deep Learning Models, Algorithms, Security & Deployment
%  Compile with: pdflatex ChemoInteract_Part2_Models_and_Algorithms.tex
% ============================================================

\documentclass[12pt,a4paper,oneside]{article}

\usepackage[T1]{fontenc}
\usepackage[utf8]{inputenc}
\usepackage{lmodern}
\usepackage[margin=2.5cm]{geometry}
\usepackage{hyperref}
\usepackage{xcolor}
\usepackage{listings}
\usepackage{longtable}
\usepackage{booktabs}
\usepackage{array}
\usepackage{amsmath}
\usepackage{amssymb}
\usepackage{tcolorbox}
\usepackage{fancyhdr}
\usepackage{titlesec}
\usepackage{enumitem}
\usepackage{microtype}
\usepackage{setspace}

\tcbuselibrary{breakable,skins}

\hypersetup{
  colorlinks=true,
  linkcolor=blue!70!black,
  urlcolor=cyan!70!black,
  pdftitle={ChemoInteract Part 2: Models and Algorithms},
  pdfauthor={Ronit},
  pdfstartview=FitH,
}

% ── Colours ──────────────────────────────────────────────────
\definecolor{codegreen}{rgb}{0,0.6,0}
\definecolor{codegray}{rgb}{0.5,0.5,0.5}
\definecolor{codepurple}{rgb}{0.58,0,0.82}
\definecolor{backcolour}{rgb}{0.97,0.97,0.97}
\definecolor{deepblue}{RGB}{0,55,120}
\definecolor{chemopurple}{RGB}{91,50,160}
\definecolor{chemogreen}{RGB}{16,140,100}
\definecolor{chemoorange}{RGB}{230,100,20}
\definecolor{noteyellow}{RGB}{255,248,200}
\definecolor{tipegreen}{RGB}{210,245,215}
\definecolor{warnorange}{RGB}{255,235,210}

% ── Listings ─────────────────────────────────────────────────
\lstdefinestyle{pythonstyle}{
  backgroundcolor=\color{backcolour},
  commentstyle=\color{codegreen}\itshape,
  keywordstyle=\color{chemopurple}\bfseries,
  numberstyle=\tiny\color{codegray},
  stringstyle=\color{chemogreen},
  basicstyle=\ttfamily\footnotesize,
  breaklines=true,
  captionpos=b,
  keepspaces=true,
  numbers=left,
  numbersep=5pt,
  showstringspaces=false,
  tabsize=4,
  frame=single,
  rulecolor=\color{gray!40},
  language=Python,
}
\lstset{style=pythonstyle}

% ── tcolorbox styles ─────────────────────────────────────────
\tcbset{
  notebox/.style={enhanced,breakable,colback=noteyellow,
    colframe=yellow!60!orange,fonttitle=\bfseries,title={Note},
    left=6pt,right=6pt,top=4pt,bottom=4pt},
  tipbox/.style={enhanced,breakable,colback=tipegreen,
    colframe=green!60!black,fonttitle=\bfseries,title={Tip},
    left=6pt,right=6pt},
  warnbox/.style={enhanced,breakable,colback=warnorange,
    colframe=orange!80!red,fonttitle=\bfseries,title={Warning},
    left=6pt,right=6pt},
}

% ── Header / Footer ──────────────────────────────────────────
\pagestyle{fancy}
\fancyhf{}
\fancyhead[L]{\small\textcolor{chemopurple}{\textbf{ChemoInteract --- Part 2}}}
\fancyhead[R]{\small\textcolor{gray}{Models, Algorithms \& Deployment}}
\fancyfoot[C]{\thepage}
\renewcommand{\headrulewidth}{0.5pt}

% ── Section styles ───────────────────────────────────────────
\titleformat{\section}{\Large\bfseries\color{deepblue}}{\thesection}{1em}{}
\titleformat{\subsection}{\large\bfseries\color{chemopurple}}{\thesubsection}{1em}{}
\titleformat{\subsubsection}{\normalsize\bfseries\color{chemogreen}}{\thesubsubsection}{1em}{}

\setstretch{1.2}
\setlength{\parskip}{6pt}

\begin{document}

% ── Title Block ──────────────────────────────────────────────
\begin{center}
{\Huge\bfseries\color{deepblue} ChemoInteract}\\[0.3cm]
{\Large\bfseries\color{chemopurple} Complete Documentation --- Part 2 of 3}\\[0.2cm]
{\large\textbf{11 Deep Learning Models, Mathematical Algorithms, Security \& Testing}}\\[0.4cm]
\begin{tcolorbox}[enhanced,width=14cm,colback=deepblue!5,colframe=deepblue,boxrule=1pt,arc=5pt]
\centering
\begin{tabular}{llll}
\textbf{Author:} & Ronit & \textbf{Version:} & 0.1.0 \\
\textbf{Framework:} & PyTorch 2.0+ & \textbf{Focus:} & Models \& Algorithms \\
\end{tabular}
\end{tcolorbox}
\end{center}

\tableofcontents
\newpage

% ================================================================
\section{11 Deep Learning Architectures}\label{sec:models}

\subsection{The Model Base Class Contract}

\begin{lstlisting}[language=Python]
class Model(nn.Module, ABC):
    @abstractmethod
    def unpack(self, batch: DrugPairBatch):
        """Extract required feature tensors from batch."""
        ...
    def forward(self, *args) -> torch.FloatTensor:
        """Execute forward pass and return probabilities in [0, 1]."""
        ...
\end{lstlisting}

\texttt{unpack()} decouples the training pipeline from model-specific input structures.

\subsection{Category 1: Fully-Connected (FC) Models}

\subsubsection{DeepSynergy (Preuer et al., 2018)}
Concatenates cell line context vector and 256-bit Morgan fingerprints of both drugs:
\[
\text{Input} = [\text{context} \,\|\, \text{left\_fp} \,\|\, \text{right\_fp}] \in \mathbb{R}^{C + 2D}
\]
Feeds through 4 linear layers with ReLU activation and Dropout (0.5) before a Sigmoid output layer.

\subsubsection{DeepDDI (Ryu et al., 2018)}
Takes concatenated drug fingerprints $[\text{left\_fp} \,\|\, \text{right\_fp}]$. Uses 9 dense layers with Batch Normalization after the first layer.

\subsubsection{MatchMaker (Kuru et al., 2021)}
Processes each drug with context using a \textbf{shared sub-network}:
\begin{enumerate}
  \item $h_L = \text{SubNetwork}([\text{left\_fp} \,\|\, \text{context}])$
  \item $h_R = \text{SubNetwork}([\text{right\_fp} \,\|\, \text{context}])$ \quad (identical weights)
  \item $\hat{y} = \text{Predictor}([h_L \,\|\, h_R])$
\end{enumerate}

\subsection{Category 2: Graph Neural Network (GNN) Models}

\subsubsection{DeepDDS (Wang et al., 2021)}
Uses dual GCN branches for molecular graphs + an MLP branch for context vector. Aggregates graph embeddings with Max Readout and combines with context embedding.

\subsubsection{DeepDrug (Cao et al., 2020)}
Uses 4 GCN layers per drug graph $\to$ Global Max Readout $\to$ Concatenation $\to$ BatchNorm $\to$ Dropout $\to$ Linear $\to$ Sigmoid.

\subsubsection{CASTER (Huang et al., 2020)}
Input is the bitwise OR union of substructures. Uses an encoder-decoder architecture with sparse dictionary coding:
\[
c = (D^T D + \lambda I)^{-1} D^T z
\]
Trained with \texttt{CASTERSupervisedLoss}.

\subsubsection{GCNBMP (Chen et al., 2020)}
Integrates **Highway Networks** ($\mathbf{y} = T \odot H(\mathbf{x}) + (1-T) \odot \mathbf{x}$), **Attention Pooling**, and **FFT-based Circular Correlation** ($a \star b = \mathcal{F}^{-1}(\overline{\mathcal{F}(a)} \odot \mathcal{F}(b))$).

\subsubsection{EPGCN-DS (Sun et al., 2020)}
Uses 2 GCN layers with Mean Readout. Combines drug representations via element-wise addition (Deep Sets principle).

\subsubsection{MR-GNN (Xu et al., 2019)}
Combines GCNs with two LSTMs: a Border LSTM tracking layer-wise feature evolution and a Middle LSTM tracking joint drug pair representations.

\subsection{Category 3: Attention Models}

\subsubsection{MHCADDI (Deac et al., 2019)}
Implements atom-to-atom cross-drug co-attention:
\[
e_{ij} = \frac{(W_K h_i)^T (W_K h_j)}{\sqrt{d}}, \quad \alpha_{ij} = \text{segment\_softmax}_j(e_{ij}), \quad m_i = \sum_{j \in B} \alpha_{ij} W_V h_j
\]
Processes drug pair graphs jointly.

\subsubsection{SSI-DDI (Nyamabo et al., 2021)}
Uses stacked GCN + LSTM blocks to model substructure-to-substructure interactions between drugs.

% ================================================================
\section{Algorithms and Mathematical Formulation}\label{sec:algorithms}

\subsection{Graph Convolution Network (GCN) Message Passing}
\[
h_v^{(l+1)} = \sigma \left( W^{(l)} \cdot \frac{1}{|\mathcal{N}(v)|} \sum_{u \in \mathcal{N}(v) \cup \{v\}} h_u^{(l)} \right)
\]
Aggregates 1-hop neighbor atom features per layer.

\subsection{Circular Correlation via FFT}
\[
(a \star b)[k] = \sum_{j=0}^{n-1} a[j] \cdot b[(j+k) \bmod n] = \mathcal{F}^{-1}\left( \overline{\mathcal{F}(a)} \odot \mathcal{F}(b) \right)
\]
Reduces correlation complexity from $O(n^2)$ to $O(n \log n)$.

\subsection{Adam Optimizer}
\[
m_t = \beta_1 m_{t-1} + (1-\beta_1)g_t, \quad v_t = \beta_2 v_{t-1} + (1-\beta_2)g_t^2
\]
\[
\theta_t = \theta_{t-1} - \frac{\alpha}{\sqrt{\hat{v}_t} + \epsilon} \hat{m}_t
\]

\subsection{Binary Cross-Entropy Loss}
\[
\mathcal{L}_{\text{BCE}} = -\frac{1}{N} \sum_{i=1}^N \left[ y_i \log \hat{y}_i + (1 - y_i) \log (1 - \hat{y}_i) \right]
\]

% ================================================================
\section{Security, Error Handling, and Troubleshooting}\label{sec:security}

\begin{longtable}{>{\bfseries}p{4cm}p{10.5cm}}
\toprule
Common Error & Cause \& Solution \\
\midrule\endhead
\texttt{AssertionError} & Invalid dataset name. Fix: use exact names (\texttt{drugcombdb}, \texttt{drugcomb}, \texttt{twosides}, \texttt{drugbankddi}). \\
\midrule
\texttt{CUDA out of memory} & Mini-batch size too large for VRAM. Fix: reduce \texttt{batch\_size} from 5120 to 512 or 256. \\
\midrule
\texttt{AttributeError: PackedGraph} & Missing GPU transfer method. Fix: import \texttt{ChemoInteract.compat} at startup. \\
\midrule
\texttt{KeyError: node\_feature} & TorchDrug version feature naming mismatch. Fix: \texttt{compat.py} resolves key mapping automatically. \\
\bottomrule
\end{longtable}

% ================================================================
\section{Deployment and Testing}\label{sec:deployment}

\subsection{Local Installation & Test Suite}

\begin{lstlisting}[language=bash]
# Clone and install in editable mode
git clone https://github.com/your-username/ChemoInteract.git
cd ChemoInteract/ChemoInteract_Main
pip install -e .

# Run pytest unit tests
python -m pytest tests/ -v
\end{lstlisting}

\subsection{Saving and Loading Models}
\begin{lstlisting}[language=Python]
results = pipeline(dataset="drugcombdb", model="deepsynergy", epochs=100)
results.save("./output_dir/")  # Saves model.pkl + results.json

# Loading model
import torch
model = torch.load("./output_dir/model.pkl")
model.eval()
\end{lstlisting}

\end{document}
